% Separation estimate for the selected finite-volume price test.
\subsection{A finite-volume separation estimate}
\label{sec:ex-lattice-trace-local-pairing}

We prove the separation estimate that converts a localized probe into a
small two-site correlation. The proof is a three-term bilinear expansion;
all quantitative inputs specific to the selected construction are written
as precise hypotheses.

Fix a finite periodic gauge configuration space $\Omega$ with normalized
Wilson expectation $\mathbb E_\beta$, and a reflection-positive cell
tiling at selected level $i$ with $n\ge1$ cells. Its admissible
observables are bounded measurable gauge-invariant functions in the
chamber-local domain. For each cell $a$, its copy $C_aF$ is the
translated pullback of $F$ under that cell's reflection; if the reflection
has odd sign parity, the copy is complex-conjugated. Define the root bound
by
\[
 \mathcal R_i(F)=\left|\mathbb E_\beta\prod_{a=1}^n C_aF\right|^{1/n}.
\]
The local positivity estimate supplies the following inequalities for
observables in the chamber-local domain:
\begin{align}
 |\mathbb E_\beta[F(\tau_{a(x)}\omega)G(\tau_{a(y)}\omega)]|
    &\le \mathcal R_i(F)\mathcal R_i(G)\quad(x\ne y),
    \label{eq:ex-lattice-root-product}\\
 \mathcal R_i(F+G)&\le\mathcal R_i(F)+\mathcal R_i(G).
    \label{eq:ex-lattice-root-sum}
\end{align}
Here $x,y$ are coarse sites in the tiling, $a(x),a(y)$ are their fine
anchor sites, and $\omega\in\Omega$ is a gauge configuration. The anchor
map is injective on the two sites under consideration. The first estimate
is the reflected-cell product bound; the second is the root-bound triangle
inequality. We use them only when the observables and all indicated sums
and differences remain in the bounded measurable, gauge-invariant,
chamber-local domain.
The local error scale at selected level $i$ is
\[
 t_i(a)=\exp\!\left(\frac{a}{10^4}(\log p_i)^{48}\right),
 \qquad p_i\longrightarrow\infty.
\]
It satisfies $t_i(a)t_i(c)=t_i(a+c)$ and is increasing in $a$ once $p_i>1$.
Since $p_i\to\infty$, also $(\log p_i)^{48}\to\infty$ and hence
$t_i(1)\to\infty$; in particular $t_i(1)\ge4$ for all sufficiently large
$i$.

\begin{proposition}[Finite-volume two-site error]
\label{prop:ex-lattice-finite-volume-pairing}
Let $P,Q,U,V$ be observables in the chamber-local domain at
selected level $i$, and let $x\ne y$ be coarse sites whose fine anchor
sites are distinct. Suppose
\begin{align}
 \mathcal R_i(P)&\le t_i(5), &\mathcal R_i(Q)&\le t_i(5),
 \label{eq:ex-lattice-pair-main-cost}\\
 \mathcal R_i(Q-V)&\le t_i(-300), &
 \mathcal R_i(P-U)&\le t_i(-200),
 \label{eq:ex-lattice-pair-local-errors}\\
 |\mathbb E_\beta[U(\tau_{a(x)} \omega)V(\tau_{a(y)} \omega)]|&\le t_i(-200).
 \label{eq:ex-lattice-pair-probe}
\end{align}
Then
\[
 |\mathbb E_\beta[P(\tau_{a(x)} \omega)Q(\tau_{a(y)} \omega)]|
 \le t_i(-100)
\]
for all sufficiently large $i$.
\end{proposition}
\begin{proof}
Write $\langle F,G\rangle_{x,y}=\mathbb E_\beta[F(\tau_{a(x)}\omega)G(\tau_{a(y)}\omega)]$. By \eqref{eq:ex-lattice-root-sum} and the assumptions,
\[
 \mathcal R_i(V)\le \mathcal R_i(Q)+\mathcal R_i(V-Q)
       \le t_i(5)+t_i(-300).
\]
The definition gives $\mathcal R_i(-F)=\mathcal R_i(F)$, since negating
$F$ changes the product of its $n$ cell copies only by the overall sign
$(-1)^n$. Hence the assumed bound for $Q-V$ also bounds $V-Q$.
Bilinearity gives the exact identity
\begin{equation}
 \langle P,Q\rangle_{x,y}
 =\langle U,V\rangle_{x,y}
  +\langle P,Q-V\rangle_{x,y}
  +\langle P-U,V\rangle_{x,y}.
 \label{eq:ex-lattice-pair-telescope}
\end{equation}
Indeed, the right hand side equals
$\langle U,V\rangle+\langle P,Q\rangle-\langle P,V\rangle
 +\langle P,V\rangle-\langle U,V\rangle$.
Using \eqref{eq:ex-lattice-root-product} on the last two terms and
\eqref{eq:ex-lattice-pair-main-cost}--\eqref{eq:ex-lattice-pair-probe}
therefore yields
\begin{align*}
 |\langle P,Q\rangle_{x,y}|
 &\le t_i(-200)+t_i(5)t_i(-300)
       +t_i(-200)\bigl(t_i(5)+t_i(-300)\bigr)\\
 &=t_i(-200)+t_i(-295)+t_i(-195)+t_i(-500)\\
 &\le4t_i(-195).
\end{align*}
The last inequality uses monotonicity in the exponent. Since $t_i(1)\ge4$,
\[
 4t_i(-195)\le t_i(1)t_i(-195)=t_i(-194)\le t_i(-100),
\]
again by monotonicity. This proves the conclusion.
\end{proof}

\subsubsection{Application to a selected response insertion}
\label{subsec:ex-lattice-mileage-application}

For the selected construction, the proposition applies with $P$ equal to
the measured grade-five marker, that is, the one-point observation of the
probe's grade-five label; $Q$ the one-point observation of the selected
grade-four insertion label, which is a separately measured grade-four
field and is not obtained from the marker; $U$ the pullback of the
localized coarse witness supplied by the probe's receipt; and $V$ the
pullback of the bounded coarse approximation supplied, at the same
finite volume and frame level, for that grade-four insertion. The
hypotheses of the proposition then take the following form: both measured
markers have
root bound at most $t_i(5)$; the local approximation satisfies
$\mathcal R_i(Q-V)\le t_i(-300)$; the probe-to-marker difference satisfies
$\mathcal R_i(P-U)\le t_i(-200)$; and the distant-site probe correlation
satisfies $|\langle U,V\rangle_{x,y}|\le t_i(-200)$. The last bound is
the probe's separated-site estimate applied to the coarse approximation
of the insertion; this is legitimate because that approximation also
satisfies the supremum bound $t_i(150000)$, the grade-four window
condition and the room condition required of a test in that estimate.
With these four receipts, the proposition proves the required two-site
price bound $t_i(-100)$.

The proof above discharges the finite-volume assembly estimate. The
coefficient-to-root and moving-depth localization arguments are supplied
below. The probe witness and its separated-site price still require the
actual Crown construction described at the end of this section. No
continuum or infinite-volume limit is used in this proposition.

\subsubsection{Production at the moving frame and localization depths}
\label{subsec:ex-trace-pairing-production}

We now supply the quantitative root and localization bounds used above.
All coefficient maps and pullbacks in this subsection belong to the one
history at $(N,p)$; the estimates are uniform in $N>k$ and $p$.
Write $j$ for the frame depth, $m$ for its first opened depth, and
$k=\operatorname{summit}(j)$. Put
\[
 x=\widehat\beta_j,\quad y=\log x,\quad H=y^{48},\qquad
 T(a)=\exp(aH/10000),\qquad
 \Delta(r,s)=\log(a_r/a_s).
\]
Thus $T$ is evaluated at the \emph{frame}, not at the summit.
The summit is the least $k>j$ with $\Delta(j,k)\ge H$.
The opened depth $m\ge j$ is the least depth with
$a_j/a_m\ge B_m^2$, where $B_m=\lceil(\log\widehat\beta_m)^2\rceil$
is the cover size. These choices are made on the common clock, before
any terminal volume or label is chosen.

Here are the clock bounds needed in the argument, including their
payment. We choose the numerical stage delays before extracting the
common history. At a fixed stage, write $r$ for its loop-path price,
$c_{\rm dir}$ for its direct quadratic-action price, and $D_G$ for the
fixed real dimension of the matrix coordinate space of the gauge group.
Let $n_{\rm plaq}$ be the number of plaquette orientations. The vacuum
majorant at coupling $z>1$ is
\[
 F_{\rm stage}(z)
 =2(1+4r+1)+1+n_{\rm plaq}(c_{\rm dir}+r^2)
       +4D_G\bigl(\log16+\tfrac12\log z\bigr)
 =C_{\rm stage}+2D_G\log z.
\]
We verify the majorization, rather than assume a bound for the treasury.
Its terms are twice the ordinary allowance plus twice the loop price,
the forced allowance, the quadratic action price, and the logarithmic
small-ball cost. The ordinary and forced allowances are at most one.
With $\vartheta$ the selected threshold, $\Lambda$ the allowed loop
cutoff, and $h=z^{-1/2}$, the loop term is
$2z(r\vartheta^3+r\vartheta^4)+\Lambda\vartheta^2$.
The clock has $\vartheta\le z^{-1/2}\exp((\log\log z)^2)$ and
$\Lambda\vartheta^2\le1$ on the late-scale tail of
Lemma~\ref{lem:rank-root-late-scale}, applied to that same clock.
The lemma's stage thresholds are among the common entry requests.
Hence $z\vartheta^3\le1$ and
$z\vartheta^4\le1$ eventually, independently of the stage constants,
and the loop price is at most $4r+1$.
The action term is exactly $n_{\rm plaq}(c_{\rm dir}+r^2)$,
since $zh^2=1$. Finally $\log(16/h)=\log16+\tfrac12\log z$.
Adding the four bounds proves $\mathcal T\le F_{\rm stage}(z)$
on a fixed coupling tail.

Now let $L$ and $c$ be the block and column constants at that stage.
For all sufficiently large $z$, impose
\[
 \log z\ge\max\{1,L,c,C_{\rm stage},2D_G+1\}.
\]
Writing $u=\log z$, we have
$C_{\rm stage}+2D_Gu\le u+2D_Gu\le u^2$.
Thus these eventual requests give, at every rung $r$ in that stage,
\[
 L_r\le\log\widehat\beta_r,\qquad
 c_r\le\log\widehat\beta_r,\qquad
 \mathcal T_r\le(\log\widehat\beta_r)^2.
\]
At each stage entry there are only finitely many active and next-stage
requests. Include the displayed bounds and their adjacent-stage versions
in the finite maximum defining the entry threshold of
Proposition~\ref{prop:trace-joint-chart-purchase}, and then take prefix
maxima to preserve earlier entries. Each request holds on an entire
coupling tail with that stage fixed, so this operation is legitimate even
though the constants vary from stage to stage. These bounds therefore
hold on the one history selected with these delays; no history is changed
after the extraction, and no uniform fixed-stage constant is asserted
across all stages.
The block lower bound $L_r\ge1024(J_{r+1}+10)$ is a choice of the
stage block before its entry threshold is paid. These purchases involve
no new finite-volume analytic estimate.

\begin{lemma}[Moving-depth bills]
\label{lem:ex-trace-pair-moving-bills}
For each fixed grade $d$, eventually
\begin{gather*}
 j\le m<k,\qquad H\le\Delta(j,k)\le2H,
 \qquad \widehat\beta_k\le2x,\\
 a_j/a_m\le y^6,\qquad
 (a_j/a_m)^4\mathcal T_m\le y^{27},\qquad
 \Delta(m,k)\ge H/2.
\end{gather*}
Put $S=\exp((a_j/a_m)^4\mathcal T_m+5)$,
$E_k=4(1+c_k)$ and $G=\exp(\Delta(m,k)/8)$.
With the fixed-grade height $\zeta_{d,r}=\widehat\beta_r^{-w_d}$,
where $w_d$ is the fixed exponent of that height schedule, these imply
\begin{gather}
 2E_k/\zeta_{d,m}\le T(100),\qquad
 2(a_j/a_m)^d S\le T(99),\qquad G\ge T(500),
 \label{eq:ex-trace-pair-localizer-bills}\\
 c_m(a_j/a_m)^d/\zeta_{d,m}\le T(2),\qquad
 eS\le T(2),\qquad 2(1+c_k)\le T(1).
 \label{eq:ex-trace-pair-root-bills}
\end{gather}
\end{lemma}
\begin{proof}
The coupling clock has increments
$\widehat\beta_s-\widehat\beta_r=\gamma\Delta(r,s)$ with fixed
$\gamma>0$. Just before the first crossing of $H$, the coupling is
at most $x+\gamma H\le2x$. The last block adds at most
$\log L_{k-1}\le\log(2x)\le H$ to the gap. Thus the gap is at
most $2H$, and $x+2\gamma H\le2x$ eventually gives the stated
summit coupling bound. In particular
$\log\widehat\beta_r\le2y$ and $B_r\le5y^2$ for $j\le r\le k$.
At $k-1$ the gap is at least $H-2y$; hence
$\exp\Delta(j,k-1)\ge(5y^2)^2$ eventually. This proves $m<k$.
Eventually $B_j>1$, so $a_j/a_j=1<B_j^2$ and $m>j$.
By minimality of $m$, the preceding ratio is less than $B_{m-1}^2$;
the last block is at most $2y$. Consequently
$a_j/a_m\le2y(5y^2)^2\le y^6$ for large $y$.
The treasury bill is at most $y^{24}(2y)^2\le y^{27}$.
Also $\Delta(j,m)\le y^6=o(H)$, which gives the last gap bound.

Every fixed product of $c_m,c_k,(a_j/a_m)^d$ and
$\zeta_{d,m}^{-1}$ is bounded by a fixed power of $x$ and $y$.
Its logarithm is $O(y+\log y)=o(H)$. Multiplication by $S$ adds
at most $y^{27}+5=o(H)$. Thus each such product is eventually
bounded by $\exp(\epsilon H)$ for every $\epsilon>0$.
Choosing respectively $\epsilon=100/10000,99/10000,2/10000$
and $1/10000$ proves the displayed bills. Finally
$G\ge\exp(H/16)\ge T(500)$.
\end{proof}

\begin{proposition}[Root and bounded localizer from actual coefficients]
\label{prop:ex-trace-pair-coefficient-production}
Let $L$ be a grade-$d$ label in the actual finite-volume history, with
its measured field in the frame chamber. If
\[
 \|\operatorname{col}_m(L)\|_1\le a_j^d T(1),
\]
then $\mathcal R_j(O_{N,p}L)\le T(5)$ eventually.
For localization use the same-door served representation constructed
below. If
$\|\operatorname{col}_k(L)\|_1\le a_j^dT(1)$, there is a bounded
coarse field $B$ at $k$, in the window of radius
$l_d=R_d+4$, such that
\[
 \|B\|_\infty\le T(25000d+10000),\qquad
 \mathcal R_j(\mathsf B_{k,N}B-O_{N,p}L)\le T(-300).
\]
The window has room inside the frame chamber, and its pullback is
chamber-local. In particular the supremum is at most $T(150000)$
for $d\le5$.
\end{proposition}
\begin{proof}
For the first assertion use insertion weight $z=T(-3)$.
The actual exponential insertion money is
$c_m a_m^{-d}\|\operatorname{col}_m(L)\|_1$; by
\eqref{eq:ex-trace-pair-root-bills} its product with $z$ is at most
$\zeta_{d,m}T(-3+1+2)=\zeta_{d,m}$.
The opened full-census exponential estimate therefore bounds its root by
\[
 \frac{\exp((a_j/a_m)^4\mathcal T_m+6)}z
 =\frac{eS}{T(-3)}\le T(5).
\]
The extra factor $e$ accounts for the difference between the single
insertion price $6$ and the discrepancy price $5$.

For localization use the actual phase-grid construction of
Subsection~\ref{sec:actual-phase-grid-approximants} with
$w=a_m^dT(400)$. Its price condition is satisfied, since
\[
 \frac{2wa_m^{-d}E_k}{\zeta_{d,m}}
 \le T(400+100)\le G.
\]
The coefficient basis in Lemma~\ref{lem:actual-grid-basis} is within
$1/4$ of the identity. Inverting that same basis therefore expresses
$L$ as a linear combination of its terminal labels with
\[
 \sum_s|v_s|\le2\|\operatorname{col}_k(L)\|_1.
\]
We retain this actual coefficient bound; replacing it by an unspecified
operator norm times $\|L\|$ would lose the moving-depth estimate.
Each primary masked field has supremum at most
$a_k^{-d}(1+c_k)$ and discrepancy root at most $S/w$.
Linear combination, pullback linearity and the root triangle inequality
give a field $B$ with error
\[
 2a_j^dT(1)S/(a_m^dT(400))\le T(99+1-400)=T(-300).
\]
Its supremum is at most
\[
 2a_j^dT(1)a_k^{-d}(1+c_k)
 \le T(2+20000d)\le T(25000d+10000),
\]
because $a_j/a_k\le\exp(2H)=T(20000)$.
The same finite combination preserves bounded measurability, gauge
invariance and the window radius $l_d$.

For completeness, the room inequality is
$2a_k[l_d+30(J_k+5)]<a_j$.
The block bounds give $J_k\le2y$, while $a_j/a_k\ge e^H$;
hence this inequality holds eventually. More generally the same proof
works for any radius bounded by a fixed power of $\widehat\beta_k$,
since such a radius and $J_k$ are $o(e^H)$.
The pullback locality and room inclusion then put $B$ in the required
frame chamber. All constructions used the same $(N,p)$, transport and
gauge doors.
\end{proof}

\subparagraph{The actual labels.}
For a fixed raw reference label $v$ of grade $d\le5$, the tested
finite-height lift at frame $j$ has coefficient bound
\[
 \|\operatorname{col}_j(Lv)\|_1
 \le a_j^d\,2(2\mathcal M_j)^2 b^d\|v\|,
 \qquad1\le b\le\operatorname{tiltcap}(j).
\]
Indeed the tested half-accuracy bound gives an inverse cost at most two;
the frame shrink inverse costs at most $2\mathcal M_j$, the
nonconstant polynomial components in these grades have degree at most
two, and the dilation weight costs at most $b^d$.
For coefficient transport from $j$ to $r\le k$, the paid per-step
logarithmic column bill divided by the log block is at most $J_s^{-2}$.
Summing gives a log bill at most
$\sup_{s\ge j}J_s^{-2}\Delta(j,r)=o(H)$.
The logarithm of the prefactor above is also $o(H)$, because
$\log\mathcal M_j=(\log\log x)^2$ and
$\log\operatorname{tiltcap}(j)=(\log\log x)^3$.
Allocate half of $H/10000$ to each bill. Uniformly for $j\le r\le k$,
\[
 \|\operatorname{col}_r(Lv)\|_1\le a_j^dT(1).
\]
This supplies both hypotheses of
Proposition~\ref{prop:ex-trace-pair-coefficient-production} for the
selected grade-four insertion, with the same tested map at each
finite height. Finitely many fixed labels require only a finite
intersection of the resulting tails.

The measured grade-five marker uses a different argument. Its odd
axial sector has homogeneous factor two, so its spine coefficients at
$r$ equal $\vartheta_r^{-2}\operatorname{col}_r(P)$.
The actual spine transport from $k$ back to $m$ costs at most two.
Since $\vartheta_m\le1$ and
$\vartheta_k^{-2}\le\widehat\beta_k$, a summit bound
$\|\operatorname{col}_k(P)\|_1\le a_j^5\widehat\beta_k^M$
gives
\[
 \|\operatorname{col}_m(P)\|_1
 \le2a_j^5\widehat\beta_k^{M+1}\le a_j^5T(1).
\]
Here the last step follows from $\widehat\beta_k\le2x$ and
$\log x=o(H)$. Thus its root is also at most $T(5)$.
This application uses the actual odd-sector spine transport; a summit
coefficient bound for an unrelated label would not suffice.


\subparagraph{The intermediate insertion and the served-representation seam.}
There is one representation issue in the localization proof that must
not be hidden by its numerical estimates. At depth $k$ the served
negative played family is augmented by the positive primary factors.
For the finite set of cell markers $I$, its coefficient is the sum over
all subsets $A\subset I$: cells in $A$ carry the local positive primary,
and the complementary cells carry the marched negative insertion.
This is the finite product expansion, with no truncation of its higher
marker coefficients. At singleton level the sum is exactly the
exponential entry of the backward coefficient family $Z_k$; its forced profile is
unchanged, because the primary has zero forced profile.
Below $k$ use that coefficient family and the same sound gauge doors. The two
transition profiles have prices
\[
 2\tau_{r+1}a_{r+1}^{-d}D_r
       (g_{r+1}+\|Z_{r+1}.x\|_1)\le1,
 \qquad \tau_{r+1}a_r^{-d}g_r\le1,
\]
as verified in Subsection~\ref{sec:actual-phase-grid-approximants}.
The centers remain the opened migrating centers, with occupancy one;
finite linear sums preserve the prescribed finite supports. These are
precisely the numerical and geometric data for the same-door offer in
Subsection~\ref{sec:finite-menu-history-handoff}. The latter invokes the
whole polynomial push, including its flagged part, of
Subsection~\ref{sec:local-transfer}; it does not permit replacing the
augmented polynomial by its singleton part.

Here is the representation construction for this augmented family.
Let $A_N$ be the negative terminal birth polynomial, and let $P_{k,N}$
be the actual iterated push from $N$ to $k$. The terminal observables are
continuous on the compact gauge configuration space, so every coefficient
of $A_N$ is integrable. The actual Markov push preserves integrability.
The local primary is bounded and measurable. Thus every coefficient of
\[
 A_k=(P_{k,N}A_N)\exp_{\mathrm{nil}}
                  \Bigl(\sum_{v\in I}z_v F_v\Bigr)
\]
is integrable: it is a finite sum of an integrable pushed coefficient
times a bounded product of $F_v$'s. Here the exponential is in the
finite nilpotent marker algebra, and $F_v$ is the reflected copy of the
positive primary. This proves integrability for the \emph{whole}
polynomial, not only for its singleton coefficients.

The negative played insertion service supplies a representation of
$P_{k,N}A_N$ with regular entries $f_v$, forced entries $\lambda_v$,
and ordinary and flagged higher arrays $H,K$. Keep $H,K$ unchanged
and replace $f_v$ by $f_v+F_v$. This represents $A_k$ exactly:
in the commutative nilpotent algebra,
$\exp_{\mathrm{nil}}(u+v)=\exp_{\mathrm{nil}}u\,
\exp_{\mathrm{nil}}v$ for zero-constant polynomials. Consequently all
higher-array support and price bounds are unchanged, and the forced
entries are unchanged. The identities for $Z_k$ above identify the new
regular entries and the old forced entries with the backward coefficient family.
There is no fresh choice of response, door, or primary coefficient.

At the next step the actual adjacent comparison of that coefficient family supplies
the finite elementary list: its regular difference is the difference of
the two exponential entries, and its flagged difference is the
finite-support forced-profile comparison transported through the same
sound door. Their two prices are the ones displayed above. Their
single-insertion identities follow from these exact adjacent identities,
while their supports and parent centers are the migrated supports and
centers already used in the negative service. The actual sound-door
response of Subsection~\ref{sec:finite-menu-history-handoff} applies to
this prepared list; its whole-polynomial push is the transition of
$A_k$, with both higher arrays retained. The integrability just proved
permits that push. Induction from $k$ down to $m$ gives a served
representation of $P_{m,k}A_k$. The weighted local entry and forced
prices at $m$ are at most $2$ and $1$, respectively, by the backward
coefficient family bounds and the purchased height inequality. Together with action
allowance $1$, the full-census estimate therefore costs $2+1+2=5$
per cell, yielding exactly $S/w$.

Finally, coefficientwise push preserves the integral, and expansion of
$A_k$ over subsets of $I$ gives, at its full mode, the product of the
reflected differences $\mathsf B_{k,N}F_v-O_{N,p}(T_Nb_s)$.
Pullback locality and the room inequality identify these with copies in
the same frame tiling. This proves the moment identity needed for the
root bound. It supplies the intermediate representation omitted from
the generic phase-grid discussion for the primary insertions used here.

\subparagraph{Assembly of the actual Crown witness.}\label{par:trace-actual-crown}
The localizer just constructed is the grade-four test of the probe;
it is not the probe's grade-five Crown witness. The source construction
of that witness has the following explicit quantitative form. At the
same summit $k$, let $X_e$ be the actual corrected primary entry minus
$\beta_{N,p,k}$ times its climate field, and let $Y_e$ be its local
finite-difference witness. Their common window must have room.
If the finite lens construction gives each summand root at most
$T(-1900000)$ and total absolute coefficient mass at most
$\widehat\beta_k^2\le T(100000)$, the root triangle inequality gives
$\mathcal R_j(\mathsf B_{k,N}Y_e)\le T(-1800000)$.
Suppose its actual healthy-region identity gives
\[
 |X_e-Y_e+c_e|\le e^{-\operatorname{prism}(\widehat\beta_k)},
 \qquad |X_e-Y_e|\le\widehat\beta_k^M
\]
with a fixed $M$ and a polynomial-size exceptional set.
Multiply by $q=(a_j/a_k)^5\le T(100000)$ and apply the finite
signed axial average. This average is contractive for these root and
supremum bounds and annihilates the constant vector $(c_e)$.
The polynomial bill is at most $T(1000)$ eventually; the healthy error
is at most $T(-600000)$ by the prism growth. Thus the averaged
witness has root at most $T(-1700000)$, global difference at most
$T(101000)\le T(500000)$, and healthy error at most $T(-500000)$.
These are the Crown bounds used in
Proposition~\ref{prop:ex-trace-crown-buffer}.

The paired estimate and healthy identity used here are constructed
below for the actual local current and the corrected-primary terms.
The finite summation, scaling and signed-average argument above then
produces the witness to which the Crown buffer applies.

\subparagraph{Payment of the paired local lens estimate.}
The moving-depth price of each local lens can also be made explicit.
Write $t=\widehat\beta_k$, $p=\operatorname{prism}(t)$ and set
\[
 w_r=T(1000000)t^{2(r-j)},\qquad j\le r\le k.
\]
Since $\widehat\beta_{r+1}\le t$, these weights obey
$w_r\le\widehat\beta_{r+1}^{-2}w_{r+1}$, exactly the two-marker
height recurrence. The number of steps is at most $\Delta(j,k)\le2H$,
so
$\log w_k\le100H+4H\log t\le p/32$ eventually. The last inequality
follows from $p=B(t)^{40}+1100$, $B(t)=\lceil(\log t)^2\rceil$,
whereas $H\le(\log t)^{48}$. In particular $1\le w_k\le e^{p/32}$.

Here is the finite-stage local payment behind this permitted weight.
Let $q=B(t)^{60}+1100$ be the lens truncation depth, and let $l$ be
its fixed cover count. Put
$c_0=\exp((s-3)l)$, where $s$ is the stage support exponent,
$c=e^{p/2}$ and $\epsilon=e^{-p/8}$. The local lease tariff is at
most $t^3$ at a late fixed-stage entry. If $u$ is that tariff times
$e^{-q}$, then
\[
 8e^p u\le\epsilon,\qquad
 4c c_0\le e^p,\qquad
 t c_0e^2w_k^2/c\le\epsilon.
\]
Indeed $q\ge4p$ eventually, while every fixed multiple of a fixed
power of $t$ is at most $e^{p/32}$. The forced allowance also satisfies
$2e^p b\le\epsilon$ by its selected small-field payment. Consequently
the local summed activity price is at most $3\epsilon$;
$(2B(t)+1)^4(3\epsilon)\le1$ and $e^{3\epsilon}-1\le1$.
These inequalities verify the local paired-marker expense with the
actual weight $w_k$. The remaining wing tariff is the one already
paid by the local finite-difference construction, at most one.
All bounds are uniform in the torus and the actual representation
within its selected parameter band.

For each reflected cell take two marker slots with the same anchor.
This has occupancy two, independently of the number of cells; the
paid occupancy schedule is eventually at least two. The anchor
separation, after displacement to the lens center, exceeds both the
lens diameter and $64B(t)p+2B(t)$, since these are polynomial in $t$
and $a_j/a_k\ge e^H$. The local lens expansion with the expense just
verified therefore forms the paired marked polynomial. Its full-mode
coefficient is the product of the reflected lens fields. The same-door
whole-polynomial transport descends it to $m$ with the recurrence for
$w_r$ above. Its regular and forced singleton entries vanish, and its
action allowance is one per marker. The full-census bound consequently
has per-cell action cost four and weight $w_m^2$, giving
\[
 \mathcal R_j(\mathsf B_{k,N}\mathrm{lens})
 \le\frac{\exp((a_j/a_m)^4\mathcal T_m+4)}{w_m^2}
 \le\exp(y^{27}+4-200H)
 \le T(-1900000).
\]
The division by a square is essential: the construction uses two
markers per cell. Thus neither the moving-depth choice nor the large
negative exponent is an additional lens smallness assumption. The local lens expansion, its wing tariff and the healthy identity
are derived in the following subsection.

\subsubsection{The local finite-difference witness and its healthy identity}
\label{subsec:trace-pair-healthy-construction}

We describe the actual fields in the Crown construction, including the
sign and the finite approximation used there. Write $h=\operatorname{peace}(t)>0$
for the finite-difference step and $Y_e$ for the localized geometric
current in direction $e$. Its support is the inner box $\mathcal B$.
For a link $a$ in this box and a color direction obtained by conjugating
a unit color by $g\in G$, let $U_{a,g,s}$ denote multiplication of that
link by the Lie exponential of $s$ times the color. Set
\[
 J_a(U,g)=\mathbf1_{\{U\text{ healthy at }a\text{ at threshold }\vartheta/4\}}
              \langle Y_e(U)_a,\operatorname{Ad}_g v_a\rangle.
\]
This is the cutoff used in the actual witness. It is bounded measurable,
local in the union of the current support and the inspected links at
$a$, and gauge covariant. Gauge covariance follows by changing $g$ by
the gauge transformation at the starting vertex; Haar measure is unchanged.
For a scalar action $F$ define
\[
 \mathcal L_{a,s}(F)(U)=\int_G
 \left[J_a(U_{a,g,s},g)e^{F(U_{a,g,s})-F(U)}-J_a(U,g)\right]dg.
\]
The actual finite-difference witness is
\[
 W_e(U)=h^{-1}\sum_{a,\nu}\omega_\nu
                  \mathcal L_{a,\lambda_\nu h}(F_q)(U),
\]
where the finite nodes $\lambda_\nu$ and weights $\omega_\nu$ are
those of the selected order-$m$ derivative formula, and $F_q$ is the
actual local action with support-cover cutoff $q$.

The change of a link stays within the local healthy chart. Indeed smoothness
of the Lie exponential on the unit ball gives
$\operatorname{chord}(\exp v)\le C_G\|v\|$ there. For the finitely
many nodes take $Z=C_G\max_\nu|\lambda_\nu|$. Since
$h=\vartheta t^{-1/24}$, eventually $64Zh\le\vartheta$.
Thus on the support of the cutoff the moved configurations have the
same large-field flags, and the current and its inspected links remain
inside the prescribed enlarged box. This is why one may use the same
local action and the same local response in every term.

To see the precise local expansion, write $F_q$ as its constant plus
its support activities $f_P$ with cover at most $q$. Activities whose
support misses $a$ have zero difference under the link change. Hence
on the support of $J_a(U_{a,g,s},g)$,
\[
 e^{F_q(U_{a,g,s})-F_q(U)}
 =\prod_{P\ni a}\left(1+
       e^{f_P(U_{a,g,s})-f_P(U)}-1\right).
\]
Expanding this product gives the empty-set term and every nonempty
cluster of changed activities, without discarding any term. The empty
term is the local directional finite difference; the other terms are
exactly the finite response corrections used to construct the corrected
primary. Locality follows by taking the union of the supports in each
term. The support-weighted sum of their absolute prices is bounded by
the exponential of the sum of the individual difference prices. This
is the elementary inequality
$\sum_{A\ne\varnothing}\prod_{P\in A}b_P
=\prod_P(1+b_P)-1\le e^{\sum_P b_P}-1$ for $b_P\ge0$.
It also proves convergence and preserves the supports when the finite
activity cutoff is enlarged.

The individual action-difference prices follow by the fundamental theorem
of calculus along the changed link. In the actual selected parameter
band their weighted sum is bounded, up to fixed stage constants and
subpower geometric factors, by
\[
 t\vartheta\,Zh+t(\vartheta^3+\vartheta^4)
            +\vartheta^2\operatorname{humble}(t)
            +\operatorname{hue}(t).
\]
Here $\vartheta=t^{-1/2}\mathcal M(t)$, $h=\vartheta t^{-1/24}$,
and the monitor and the displayed geometric factors are subpowers of
$t$. The first term is $Z\mathcal M(t)^2t^{-1/24}$;
the second and third have strictly negative powers as well. Thus this
sum tends to zero, and is eventually at most one. This is the wing
tariff used in the paired-lens payment, now for these actual link changes.
The bound is uniform in the torus because the price is a support-weighted
sum based at $a$, not a sum over all translated cells.
The cutoff current is at most $t$ eventually: its derivative bound is a
fixed subpower current bound. The same expansion bounds each lens by
$t(e^2+1)$, and therefore the finite weighted sum defining $W_e$ is
at most $t^2$ after paying its fixed node weights, the subpower number
of links and $h^{-1}\le t^{2/3}$ eventually. These are the global bounds required above.

Choose a finite activity cabinet $F_{q,K}$ using the very same activities
as $F_q$. Their uniformly summable prices imply uniform convergence of
the action differences and hence of their exponentials. More explicitly,
if the omitted absolute difference price is $\eta_K\to0$ and all full
difference prices are at most $C$, then the difference of exponentials
is at most $e^C(e^{\eta_K}-1)$, uniformly in $U$ and $g$.
Multiplication by the bounded cutoff current and integration over Haar
measure preserve this convergence. If
$D_m=\sum_{a,\nu}|\omega_\nu|$, choose the common cabinet sufficiently
large that each lens error is at most
$h e^{-p}/(1+D_m)$. Then
\[
 \left|W_e-h^{-1}\sum_{a,\nu}\omega_\nu
              \mathcal L_{a,\lambda_\nu h}(F_{q,K})\right|
 \le e^{-p}.
\]
There are finitely many directions, so one cabinet suffices for all of
them. Enlarging it to contain the finite action and loop cutoffs used by
the corrected primary preserves these tail bounds. The cabinet is
chosen at the given finite volume; no volume-uniform bound on its number
of terms is needed.

On the common healthy region the cutoff is one for every moved
configuration. For the healthy identity we refactor the finite cabinet,
keeping its short main and loop action in the core. Put
$F=-\beta M-L$, where $M$ and $L$ are the actual finite main-action
and loop-action sums. The exact local action identity is
\[
 F_{q,K}(U_{a,g,s})-F_{q,K}(U)
 =F(U_{a,g,s})-F(U)+\sum_{P\in\mathcal A_a}\delta_P(U,g,s),
\]
where $\mathcal A_a$ contains only the residual changed activities,
excluding those already placed in $F$. Thus its exponential factors as
$e^{F(U_{a,g,s})-F(U)}\prod_{P\in\mathcal A_a}(1+(e^{\delta_P}-1))$.
This differs from the expansion of the whole action used for the lens
price: here the core exponential stays outside the product. Let $D_Y$
denote the current directional derivative. The empty residual cluster is
\[
 h^{-1}\mathcal D_m(F,Y)=\operatorname{div}Y+D_YF+K_m(F,Y),
\]
where $K_m$ is the exact finite-difference remainder (defined by this
identity). Write $R_{a,A}$ for the nonempty-cluster corrections from the
product above and put
$\mathcal C=h^{-1}\sum_{a,\nu,A\ne\varnothing}\omega_\nu R_{a,A}$.
The corrected primary is precisely the sum
\[
 Z_e=\beta(B_e-D_YM)
       +(\operatorname{div}Y_e-\operatorname{div}Y_e(1))
       -D_YL+K_m(F,Y_e)+\mathcal C.
\]
These are the actual residual, centered divergence, loop, finite-difference
and nonempty-cluster terms; each is constructed with the same current,
finite cabinet and finite action sums. Combining the empty and nonempty
terms and using $D_YF=-\beta D_YM-D_YL$ gives the exact identity
\[
 h^{-1}\sum_{a,\nu}\omega_\nu
       \mathcal L_{a,\lambda_\nu h}(F_{q,K})
       =Z_e-\beta B_e+\operatorname{div}Y_e(1).
\]
Thus for $X_e=Z_e-\beta B_e$ and
$c_e=\operatorname{div}Y_e(1)$ the cabinet estimate proves
$|X_e-W_e+c_e|\le e^{-p}$ on the common healthy region.
This supplies the healthy identity with its actual constant and sign.
Applying the finite signed axial average kills $(c_e)$, exactly as in
the Crown assembly above.

\subparagraph{Selecting and localizing the actual geometric current.}
The current in this construction is selected from the inherited geometric
fields on the same finite-volume history. The selection uses the four-jet
bounds and the finite family of quadratic tests, not compactness of all
smooth fields. At a fixed depth $k$, the inherited finite-height
construction has error bounds of the form
$2^{-(r+1)}d_j(\widehat\beta_k)$ for jet orders $2\le j\le4$.
The history used here is the one bound to its own transfers and doors
by Proposition~\ref{prop:ex-lattice-selected-course-binding}. We now
include the geometric purchases in that history's finite-stage ledger.
Fix a stage profile, kernel choice, and required differentiation order.
The local adjacent-field estimates give positive allowable analytic
radius $a_*$, a carrier constant, and an early derivative constant
$H_*$. Choose the integer block large enough for the inequalities
below, which involve $H_*$ only; the late derivative estimate then
gives the finite constant $D_*$ and radius $b_*$. With that block fixed, the coupling increment
$v=\gamma\log L$ is fixed. Increasing the stage entry ensures
\[
 (\log t)^{-1}\le\min(b_*,a_*/2),\quad
 (\log(t+v))^{-1}\le c_*,\quad
 \frac{\log(t+v)}{\log t}\le L,\quad
 \lceil(\log t)^2\rceil\ge R_*.
\]
Here $b_*,c_*>0$ are the late derivative and canonical-top radii,
and $R_*$ is the finite geometric support requirement. The canonical
path-field jet bill is also bounded by its logarithmic budget after
increasing this same entry. Each inequality follows respectively from
inverse-log decay, the log-ratio limit one, cover divergence, or a fixed
jet bill divided by a positive log power. The bank inequalities below
are paid at that entry as well. The maximum over the finitely many
active and next-stage profiles is an admissible common-calendar
purchase. Consequently every sufficiently late adjacent transition of
this very history has the required geometric derivative, support and
bank bounds. No separate history with better estimates is substituted.

Here the inheritance is the geometric pullback through the actual
selected sound door, followed by the fixed main, flux, leaf and trace
averages; it is not the coefficient chart's half-accuracy map.
Concretely a scalar component of engineering degree $d$ is pulled back
through the selected stationary interpolation and averaged with its
specified finite sheaf, then multiplied by $L^d$. The vector current
uses the differentiated interpolation and factor $L^5$. The stress
uses the flux sheaf, the current and error use the main sheaf, and the
auxiliary scalar uses the leaf sheaf plus the trace correction of the
stress. The twist uses its leaf term and the two antisymmetric trace
corrections. These are the exact discrete product-rule corrections:
they preserve the trace, symmetry and divergence identities at every
step. The interpolation, its derivative and the path entering these
formulas are the ones of the actual selected transition, as bound in
Proposition~\ref{prop:ex-lattice-selected-course-binding}; there is
no additional choice of a free interpolating field. Its
finite-height estimates are obtained as follows. At the top finite
height use the following explicit fields. Let
$F_{ab}(z)=\operatorname{curv}(U(\square_{ab}(z)))$ be the oriented
plaquette curvature, with the curvature map and color inner product
of the fixed gauge group. Set
\[
 J^0_{ab}(z)=\frac{\delta_{ab}}4\sum_{i,l}
          \langle F_{il}(z),F_{il}(z)\rangle
        -\sum_i\langle F_{ai}(z),F_{bi}(z)\rangle.
\]
The current $Y^0_b(z)$ is supported on links starting at $z$ and has
color component $F_{ba}(z)$ on the link in direction $a$. Put
$Q^0=0$, $\mathrm{twist}^0=0$, and
$E^0_b=\sum_a\nabla_a J^0_{ab}-\mathcal A(Y^0_b)$, where
$\mathcal A$ is the actual path action's directional variation.
Trace zero and symmetry of $J^0$ follow by summing its displayed
formula and symmetry of the inner product; the divergence identity
holds by the definition of $E^0$. Gauge covariance follows from
conjugation covariance of plaquette curvature. At the flat field
$F_{ab}=0$, so the scalar fields have zero constant and first jets.
Their four-jet bounds follow by differentiating the finite plaquette
products and the fixed smooth curvature map, together with the actual
path-action derivative bounds. The quadratic tests are obtained by
replacing each plaquette curvature by its first derivative, the discrete
curl. On the fixed polynomial test jets this curl agrees with the free
curvature as soon as the period contains the test stencil without wrap;
products give the stress quadratic form and pairing the current's
linear curl with the test direction gives its mixed grade-five form.
Thus these tests are eventually exact on the growing-period tail.
Their
weighted jets through order four have fixed bounded prices, their
quadratic tests converge to the prescribed free forms as the period
grows, and their geometric identities hold exactly. At each adjacent
step apply the actual inherited derivative estimate, reserving one
third of the next jet budget for each summand in the stress, auxiliary
and twist formulas. If $A_*$ is the fixed total averaging debt and
$H_*$ the derivative constant, choose the stage block $L$ so that
\[
 3A_*2^{13}H_*^2/L<1/2,\qquad
 3L^5A_*(2H_*/L^2)^s<1/2\quad(s=3,4).
\]
Such an integer block exists because the powers of $L$ are $-1,-1,-3$.
The second-jet principal contribution is bounded by
$A_*2^{13}H_*^2(C+(d_2(t')+C)/L)$; dividing by $d_2(t)$,
with $d_s(t)=(\log t)^s$ and $t'=t+\gamma\log L$, its limit is
$A_*2^{13}H_*^2/L<1/6$.
For orders three and four the full majorant, before division by the
coarse budget, is
\[
 L^5A_*\left[(2H_*/L^2)^s d_s(t')+
   \sum_{\pi\ne\mathrm{atomic}}D_*^{|\pi|}d_{|\pi|}(t')\right].
\]
The sum is over the finitely many non-atomic ordered derivative
partitions, each with $|\pi|<s$; it is retained in the estimate.
The leading quotient tends to
$L^5A_*(2H_*/L^2)^s<1/6$. Every other chain-rule term has
strictly fewer than $s$ derivative factors, so its quotient contains
$(\log t)^{r-s}$ with $r<s$ and tends to zero. Thus a late stage
entry pays all three summands and halves the error budget at each
step. The same exact pullback and averaging identities preserve the
quadratic tests and the geometric field equations. Induction over the
finite number of adjacent steps gives the displayed
$2^{-(r+1)}d_j$ error bounds. Only after fixing that number do we
pass to the large terminal-depth and large-period tail of the same
history. These are the stage block and entry purchases already allowed
by the common selected-history calendar.

Given a positive error $\epsilon$, first choose a finite $r$ making
these three bounds less than $\epsilon$. The fixed-height convergence
then gives, on the common $(N,p)$ tail, the inherited fields with jet
price at most $d_j(\widehat\beta_k)+1$ and quadratic-test error at
most $\epsilon$. These statements hold uniformly over sites: otherwise
choose a violating site at each bad volume, contradicting the same
fixed-height estimate for that sequence of sites. There are only
finitely many component and quadratic-test indices, so one tail suffices.
Take $\epsilon=\operatorname{shock}(\widehat\beta_k)$ at this fixed
$k$, and average over all translations on that same torus.
The inherited field equations are linear, the price bounds convex,
and the quadratic tests translate with their base point, so this average retains all bounds
and gives the required exact symmetry of the geometric family.
In particular the actual current has the common exponentially weighted
four-jet price $\operatorname{trophy}(t)$ at decay rate
$a=\operatorname{east}(t)$.

Localize this current at radius $r=\operatorname{rim}(t)$ using the
same smooth local extension as the other geometric fields. The
localization estimate in the charge construction bounds its lost
four-jet price by
\[
 C\,\operatorname{momentPrice}_{P}(a/4)
                 e^{-ar/4}\operatorname{trophy}(t).
\]
The retained current is supported in the box of radius $4(r+2)$.
For any fixed derivative order $s$ its composition derivative price is
at most $P_s(r)\operatorname{trophy}(t)$, with a fixed polynomial
$P_s$. To obtain this last statement, differentiate the local extension:
at each derivative order there are finitely many product and chain-rule
terms, each involving a bounded derivative of the fixed extension and
at most a fixed power of the number of sites in the radius-$r$ box.
The four-jet coefficient price supplies the coefficients; the extension
has only these four polynomial jets before composition. Thus neither
this polynomial nor its constants depend on the ambient torus.
The clock choices give
\[
 C\operatorname{momentPrice}_{P}(a/4)e^{-ar/4}
                    \operatorname{trophy}(t)\le t^{-70},\qquad
 P_s(r)\operatorname{trophy}(t)\le\mathcal M(t)^{1/200}
\]
eventually, for each fixed $s$. The first is the exponential-tail
purchase; the second follows by taking logarithms of the fixed
polynomial bill against the monitor. These are the local current and
derivative bounds used in the cutoff and finite-difference construction.
All fixed orders required at an active stage are paid together before
that stage is used. This fixes the current itself before the finite
cabinet approximation and before the corrected-primary coefficients
are chosen.


The same localization is applied to the scalar geometric stress, its
auxiliary scalar field and its first-difference fields. Their local
extensions use the same cutoff; the finite algebraic identities between
them are preserved by taking the prescribed shifted linear combinations.
The product and chain-rule argument just given bounds each fixed-order
jet by a polynomial in the localization radius times its inherited
price. Passing to a first difference adds only finitely many translated
copies and one fixed enlargement of the support. The geometric current,
stress, auxiliary field, twist fields and divergence therefore have
common fixed-order analytic bounds at most $\mathcal M(t)^{1/50}$
after increasing the stage entry. Their zero and first scalar jets are
preserved by the local extension. The healthy quadratic Taylor
conversion and the resulting reference-field price in the charge
construction consequently apply to this complete family, including its
first differences; they are not inferred from a bound on the current
alone.
